\section{\ref{sec:debugging}장. 기계 디버깅}

전극이 무엇을 듣는지, 활동의 저차원성, 딱딱한 배열을 쓴 인터페이스의 성능, 뇌와 디코더의 공동 학습, 자극으로 생기는 시각의 점은 동료 심사를 거친 실험에서 측정되었다. 데이터 흐름의 추정은 본문의 산술이다. 사람에게 이식한 Neuralink 장치에 대해서는 확인한 범위에서 주로 회사 자체가 데이터를 발표했다.

{\footnotesize
\setlength{\tabcolsep}{3pt}\renewcommand{\arraystretch}{1.15}
\begin{longtable}{>{\raggedright}p{0.5cm}>{\raggedright}p{4.0cm}>{\raggedright}p{5.6cm}>{\raggedright}p{2.3cm}>{\raggedright\arraybackslash}p{1.7cm}}
\toprule
No. & 명제 & 실험과 측정한 것 & 출처 & 상태 \\
\midrule\endfirsthead
\toprule
No. & 명제 & 실험과 측정한 것 & 출처 & 상태 \\
\midrule\endhead
1 & 전극은 반경 약 100마이크로미터 안의 뉴런, 보통 한 개에서 몇 개의 스파이크를 듣고, 모양으로 구별한다 & 쥐, CA1 영역, 세포 안과 세포 밖 동시 기록: 세포 밖 스파이크의 모양은 세포 안 스파이크의 폭과 진폭을 따른다. 추정: 테트로드는 뉴런 $60$--$100$개를 구별할 수 있지만 실제로는 약 여섯 개를 골라낸다. & \cite{Henze2000extra} & 측정됨 \\
2 & 뇌에는 뉴런 $8.6\cdot10^{10}$개가 있고, 프로브는 대략 1억분의 1을 듣는다 & 녹인 조직의 현탁액에서 핵을 셈: 사람 뇌의 뉴런 $86\cdot10^9$개. 비율은 본문의 산술. & \cite{Azevedo2009} & 측정됨 \\
3 & 운동 피질 뉴런 수백 개의 활동은 공통 변수 열 개 남짓에 들어간다 & 원숭이 운동 피질 기록의 리뷰: 집단 활동은 많은 뉴런이 공유하는 소수의 잠재 변수로 기술된다. & \cite{Gallego2017} & 측정됨 \\
4 & 이웃의 잡음은 상관되어 있고, 뉴런 100개는 1000개와 거의 같은 양을 전달한다(명제~\ref{prop:pooling}) & 원숭이 시각 피질: 이웃 뉴런의 잡음 상관 계수는 약 $0.1$이며, 평균의 이득은 뉴런 100개에서 포화된다. 정보를 제한하는 상관의 이론. & \cite{Zohary1994, MorenoBote2014} & 측정됨 \\
5 & 원시 흐름 $2\cdot10^8$~bit/s 대 스파이크 $8\cdot10^4$~bit/s~\eqref{eq:probedata} & 스파이크 흐름의 용량~\eqref{eq:spikeentropy}에 따른 본문의 산술. 디지털화 매개변수는 실제 값이다: Neuralink 칩은 채널당 $19.3$~kHz, $10$비트. & 이 장의 유도; \cite{Rieke1997, Musk2019} & 이론 \\
6 & Neuralink의 첫 시스템: 칩이 증폭하고 디지털화하며, 원시 흐름은 케이블로 나가고, 스파이크는 외부 프로그램이 골라낸다. 칩 위의 스파이크 검출, 밀폐된 본체, 무선 통신과 무선 전력은 회사 발표에 따르면 사람용 장치에 있다 & 회사 논문: 원시 흐름 전체가 USB-C 케이블로 나가고, 스파이크는 칩 밖의 프로그램이 실시간으로 골라낸다. 탑재 압축, 무선 전력과 전송은 임상 장치의 계획으로 언급되었다. 사람에게 이식한 장치는 회사 발표뿐이다. 칩 위에서 스파이크를 골라내야 한다는 것은 \eqref{eq:probedata}에서 나온 본문의 결론이다. & \cite{Musk2019}; Neuralink 보고, 동료 심사 논문 없음 & 측정됨(회사 논문); 사람은 회사 보고 \\
7 & 실 $96$가닥에 전극 최대 $3072$개, 실마다 $32$개; 실은 로봇이 혈관을 피해 꽂는다 & 회사 논문: 실 $96$가닥에 전극 $3072$개, 실마다 $32$개; 로봇이 1분에 실 $6$가닥(전극 $192$개)을 꽂는다; 배열을 오래 이식한 쥐에서 전극의 최대 $70\,\%$가 스파이크를 듣는다. & \cite{Musk2019} & 측정됨(회사 논문) \\
8 & 사람에서는 실 $64$가닥에 약 1000채널; 실 일부가 밀려났다; 커서는 약 $10$~bit/s & 회사 발표뿐. PubMed 검색에서 사람 대상 결과를 담은 동료 심사 논문을 찾지 못했다. 본문의 단서와 일치한다. & Neuralink 보고; 동료 심사 논문 없음 & 측정됨(회사 보고) \\
9 & 전극 100개 남짓의 배열을 쓴 인터페이스가 커서를 조작한다 & 마비 환자, 일차 운동 피질에 전극 $96$개: 손상 3년 뒤에도 팔을 움직이려는 의도가 스파이크를 바꾼다. 디코더가 “신경 커서”(메일, TV), 의수 손과 로봇 팔 조작을 가능하게 한다. & \cite{Hochberg2006} & 측정됨 \\
10 & “생각 속의 손”으로 글쓰기: 분당 약 $90$자, $8$~bit/s 미만 & 손이 마비된 사람, 손으로 쓰려는 시도, 순환 신경망 디코더: 실시간으로 분당 $90$자, 정확도 $94.1\,\%$, 자동 교정 뒤 $99\,\%$ 이상. 비트 추정은 본문의 산술. & \cite{Willett2021} & 측정됨 \\
11 & 말하려는 시도를 분당 약 $60$단어로 글로 바꾼다 & 또렷한 말을 잃은 사람, 피질의 배열: 분당 $62$단어; 어휘 $50$단어에서 단어 오류 $9.1\,\%$, 어휘 $125\,000$단어에서 $23.8\,\%$. & \cite{Willett2023} & 측정됨 \\
12 & 인터페이스는 용량의 작은 일부만 꺼낸다: 뉴런 하나는 초당 수십 비트, 100개는 수천 비트(명제~\ref{prop:probe}) & 상한~\eqref{eq:spikeentropy}은 이론; 8행, 10행과의 비교는 본문의 결론. 병목이 전선이 아니라 저차원성과 부호에 대한 무지라는 것은 직접 실험으로 확인되지 않았다. & \cite{Rieke1997} & 이론 \\
13 & 뇌와 디코더는 서로에게서 배운다 & 원숭이가 커서를 조작; 디코더가 닫힌 루프에서 맞춰지는 동안 뉴런도 다시 학습한다: 정확도가 오르고, 기술이 유지되며, 자기 팔의 움직임에 덜 방해받는다. 학습 속도를 떼어 놓으라는 조언은 공학 규칙이며, 본문에 별도 실험은 없다. & \cite{Orsborn2014coad} & 측정됨 \\
14 & 전극을 통한 전류는 주변 뉴런과 전선을 유형을 가리지 않고 흥분시킨다 & 생쥐와 고양이 피질, 2광자 칼슘 기록: 약한 전류조차 수십 마이크로미터 지름의 부피에서 축삭을 직접 흥분시켜, 최대 수 밀리미터 떨어진 드문드문한 뉴런을 흥분시킨다. 즉 흥분은 비선택적이지만 빽빽하지는 않다. & \cite{Histed2009stim} & 측정됨 \\
15 & 시각 피질 자극은 점을 켠다. 최대 $16$개 점을 동시에 켜면 몇몇 글자와 물체의 경계를 알아볼 수 있다 & 시각 장애인, 시각 피질에 전극 $96$개를 $6$개월 이식: 전극 하나의 문턱값 $66.8\pm36.5$~\textmu A; 무리를 동시에 자극하면(최대 $16$전극, 자극기의 한계) 문턱값이 낮아지고 구별 가능한 무늬가 생기며, 몇몇 글자와 물체의 경계를 알아본다. & \cite{Fernandez2021} & 측정됨 \\
16 & Neuralink의 두 번째 방향은 시각 피질에 신호를 넣는 장치다 & 개발 중이라는 회사 발표뿐. 작동에 관한 동료 심사 데이터는 찾지 못했다. & Neuralink 보고 & 가설 \\
17 & 브로드캐스트 신경전달물질의 농도는 조직 안의 화학 센서로 잴 수 있다 & 쥐 뇌 절편, 탄소 섬유 고속 순환 전압전류법: 세로토닌의 방출과 재흡수를 측정; 물질이 시냅스 밖으로 나간다. & \cite{Bunin1998} & 측정됨 \\
\bottomrule
\end{longtable}}
